\CUSTOMstruct{ВВЕДЕНИЕ}

Выпускная квалификационная работа является обязательным этапом государственной итоговой аттестации. 
Она представляет из себя некоторую работу, которая показывает навыки и уровнь подготовки студента, сформированные за время обучения.
Подготовка к написанию выпускной квалификационной работы требует навыки работы с проектной и технической документацией. 

Данный отчёт отражает результаты подготовительных этапов к написанию выпускной квалификационной работы. 

Цель практики заключалась в подготовке к написанию ВКР, путём осознания методологий составления проектных документов, изучению требований и применению на практике знаний.

Чтобы выполнить поставленную цель были проделаны следующие этапы, приведённые в таблице~\ref{tab:practice-stages}.

\begin{table}[h]
  \centering
  \caption{Этапы и задания практики}
  \label{tab:practice-stages}
  \renewcommand{\arraystretch}{1.2}
  \begin{tabularx}{\textwidth}{|>{\raggedright\arraybackslash}p{0.28\textwidth}|>{\raggedright\arraybackslash}X|}
    \hline
    \textbf{Название этапа} & \textbf{Задание} \\
    \hline
    Инструктаж обучающегося & Инструктаж обучающегося по ознакомлению с требованиями охраны труда, техники безопасности, пожарной безопасности, а также правилами внутреннего трудового распорядка \\
    \hline
    Лекции & Посетить 2 занятия. \\
    \hline
    Шаблон для ВКР & Подготовить шаблон для ВКР по требованиям университета. Создать шаблоны шрифтов, задать структуру ВКР. Прислать сформированный на основе шаблона текстовый документ в формате doc или txt, или подготовленный шаблон в LaTeX или Typst. Задание выполняется полностью самостоятельно, без использования заготовок. Сформировать документ или архив и прислать на проверку на электронную почту преподавателя. \\
    \hline
  \end{tabularx}
\end{table}

\newpage
\begin{table}[h]
  \centering
  \caption*{Продолжение таблицы~\ref{tab:practice-stages}}
  \renewcommand{\arraystretch}{1.2}
  \begin{tabularx}{\textwidth}{|>{\raggedright\arraybackslash}p{0.28\textwidth}|>{\raggedright\arraybackslash}X|}
    \hline
    \textbf{Название этапа} & \textbf{Задание} \\
    \hline
    Техническое задание по ВКР & Написать техническое задание по ВКР. Техническое задание должно содержать: наименование, назначение, основание для разработки, функции, структура, пользовательский интерфейс, требования к надёжности, требования к безопасности, условия эксплуатации и проч.\ важные требования, документация, стадии и этапы разработки, порядок контроля и приёмка. В случае отсутствия какого-либо пункта необходимо обосновать его отсутствие. Сформировать в виде документа в формате pdf и прислать на проверку на электронную почту преподавателя. \\
    \hline
    Оформление отчётных документов в соответствии с требованиями & 1.~Должно быть подробное описание выполнения задач по этапам. Результаты задания необходимо разместить в приложения. 2.~Оформление отчёта должно быть выполнено в соответствии с методическим пособием (\url{https://books.ifmo.ru/file/pdf/2622.pdf}). 3.~Структура документа: титульный лист, введение, основная часть, заключение, приложения. 4.~В основной части подробно описывается выполнение задач 2--4 этапов, в приложении помещаются результаты данных этапов. 5.~Отчёт необходимо подгрузить в модуле практика как <<письменный отчёт>>. \\
    \hline
    Получение отзыва руководителя практики & Получить отзыв с оценкой у руководителя практики в модуле Практики \\
    \hline
  \end{tabularx}
\end{table}